% analysis_section.tex (2026-10-18): short analysis paragraphs for the paper, condensed from distilc_theory.tex
% (full statements and proofs -> appendix). Citations: analysis_refs.bib.

\paragraph{Assumptions about goal continuity.}
Let $g'$ be an unflown goal at distance $d$ from a flown goal $g$, with both rollouts visiting the same sequence of
contact events. We assume that, within this mode sequence, the real error and the sensitivities are Lipschitz in the
goal, $\|\mbf{e}_{g'}-\mbf{e}_g\|\le L_e d$ and $\|\mbf{G}_{g'}-\mbf{G}_g\|\le L_G d$. This follows from continuous
dependence on parameters for smooth dynamics \citep[Ch.~3]{khalil2002nonlinear} and from piecewise differentiability
for hybrid systems with transversal events \citep{burden2016pwd}, and it fails across mode boundaries (e.g.\ flat
vs.\ box landings). It holds when the sim-to-real discrepancy is robot-wide (gain, mass, friction, delay) rather than
specific to one goal.

\paragraph{Assumptions about initial state perturbation and basins of convergence.}
Initial states $\mbf{x}_0\sim p_0(\cdot\mid g)$ and sensing enter the measured error as zero-mean noise of bounded
variance. Unlike classical ILC, which needs repeated initial conditions \citep{heinzinger1992stability}, the policy's
feedback rejects part of each perturbation within the rollout. The analysis is local: it needs a neighbourhood of the
current policy in which actions do not saturate and the mode sequence is fixed, and an initial error small enough to
lie in the basin of convergence. Placing the policy in that basin is the role of the simulation stage. Rollouts that
change contact sequence or fall leave it; their errors are masked and worse updates are rolled back.

\paragraph{Proof of convergence under assumptions.}
A parameter change acts on the closed loop like an additive correction, so to first order the error evolves as
$\mbf{e}^{(i+1)}=(\mbf{I}-\beta\mbf{M})\mbf{e}^{(i)}$, with $\mbf{M}=\mbf{G}_{\mathrm{real}}\mbf{S}\hat{\mbf{L}}$:
the step requested by the nominal model,
$\hat{\mbf{L}}=(\mbf{G}^\top\mbf{Q}\mbf{G}+\lambda\mbf{I})^{-1}\mbf{G}^\top\mbf{Q}$, realized by the network refit
$\mbf{S}$, and passed through the true closed loop $\mbf{G}_{\mathrm{real}}$. We assume \emph{positivity},
$\mbf{e}^\top\mbf{Q}\mbf{M}\mbf{e}\ge\mu\|\mbf{e}\|_{\mbf{Q}}^2$ and $\|\mbf{M}\|_{\mbf{Q}}\le\Gamma$: the nominal
correction reduces the real error to first order. It holds, for example, for multiplicative model errors of any size,
$\mbf{G}_{\mathrm{real}}=\mbf{P}\mbf{G}$ with $\mbf{P}+\mbf{P}^\top\succ0$ \citep{harte2005inverse}, and it is the
approximate-model ILC condition \citep{bristow2006survey}. Then, for $\beta\le\mu/\Gamma^2$,
$\|(\mbf{I}-\beta\mbf{M})\mbf{e}\|_{\mbf{Q}}^2\le(1-2\beta\mu+\beta^2\Gamma^2)\|\mbf{e}\|_{\mbf{Q}}^2\le(1-\beta\mu)\|\mbf{e}\|_{\mbf{Q}}^2$,
and with the Taylor remainder bounded on the neighbourhood the error converges linearly. With trial noise and
decreasing step sizes, convergence is almost sure by the Robbins--Siegmund theorem \citep{robbins1971convergence}.
Moreover, if the network can realize any correction and $\mbf{G}$ has full row rank, the update stops only when
$\hat{\mbf{L}}\mbf{e}=\mbf{0}$, i.e.\ $\mbf{e}=\mbf{0}$: the fixed points are zeros of the \emph{real} error,
whatever the gap $\hat f\neq f$. The model error affects the rate, not the limit (Appendix~\ref{app:theory}).

\paragraph{Why we don't use fitted value functions.}
Deep off-policy ILC \citep{gurumurthy2023deepilc} fits a critic to a value-gradient loss propagated through
approximate simulator Jacobians. We regress onto the ILC step instead, for three reasons. Holding the next action
fixed makes the critic's target the open-loop gradient, which for a feedback policy was uncorrelated with the exact
closed-loop gradient in our simulator ($r=0.00$), while training on closed-loop targets diverged, since the targets
depend on the actor's gains that follow the critic \citep[Sec.~11.3]{sutton2018rl}. The exact closed-loop step is
already available from $\mbf{G}$, so a learned approximation only adds error. And a fitted critic moves the fixed
points to the zeros of its own approximation, reintroducing the model bias that the measured residual removes
\citep{abbeel2006inaccurate}.

\paragraph{Why single corrections correct a neighbourhood of goals.}
To first order, the refit changes the policy at every input, not only at the targets:
$\delta\pi(\mbf{x},g)=\Theta\big((\mbf{x},g),\mcal{D}\big)\big(\Theta(\mcal{D},\mcal{D})+\rho\mbf{I}\big)^{-1}\mbf{v}^\star$,
where $\Theta$ is the empirical neural tangent kernel \citep{jacot2018ntk,lee2019wide}. A correction measured at goal
$g_j$ therefore also moves the policy at goals the kernel couples to $g_j$. Under the continuity assumptions above and
positivity at $g_j$, this movement reduces the error at an unflown goal within a radius set by $L_e$, $L_G$ and the
kernel's decay. Across mode boundaries there is no such guarantee, and convergence is guaranteed only at the flown
goals.
